\subsection{Partisan Clustering in American Politics}

Geographic partisan sorting in the United States has accelerated since the
1990s \citep{bishop2009bigSort, rodden2019why}.  \citet{chen2013unintentional}
demonstrated that even random redistricting algorithms produce a systematic
Republican seat bonus because Democratic voters are more densely concentrated
in urban centres, creating large wasted-vote surpluses in Democratic-leaning
districts.  \citet{rodden2019why} extended this analysis, arguing that the
geographic concentration of Democratic voters---the ``Rodden gap''---is the
primary source of partisan asymmetry in U.S. legislative elections, with
intentional gerrymandering as a secondary factor.

This literature suggests a puzzle: if geographic sorting already produces safe
seats, what does deliberate safe-seat gerrymandering add?  The answer is
directional: geographic sorting produces safe \emph{Democratic} districts in
urban areas, but the Republican seat bonus arises from dispersing Democratic
voters across suburban and rural Republican districts.  Intentional Republican
gerrymanders amplify this dispersion, cracking urban-adjacent Democratic
precincts and reinforcing Republican safety in suburban seats.

\subsection{The Efficiency Gap}

The Efficiency Gap (EG), introduced by \citet{stephanopoulos2015partisan},
measures the difference in ``wasted votes'' between the two parties.  A
wasted vote is any vote cast for a losing candidate (all such votes are
entirely wasted) or any vote cast for a winning candidate beyond what was
needed to win (excess votes above 50\% plus one are wasted).  Formally:

\begin{equation}
  EG = \frac{WV_D - WV_R}{TV}
\end{equation}

where $WV_D$ and $WV_R$ are total wasted votes for Democrats and Republicans
respectively, and $TV$ is total votes cast.  A positive EG favours Republicans
(Democrats waste more votes); a negative EG favours Democrats.

PSDs can be designed to minimise the absolute EG: by distributing safe seats
proportionally between the parties, the algorithm ensures neither party wastes
systematically more votes than the other.  This is in contrast to enacted
gerrymanders, which typically produce large EGs in favour of the mapmaking
party.

\subsection{Geographic Sorting vs.\ Intentional Packing}

The key distinction for legal analysis is between \emph{organic} geographic
sorting (which is constitutionally unregulated) and \emph{intentional} packing
(which may violate state constitutional provisions on partisan fairness).
PSDs help operationalise this distinction: a PSD map created from geographic-
only sorting (partisan weights but no geographic boundary preference) represents
the ``natural'' safe-seat level arising from residential patterns.  An enacted
plan with higher safe-seat counts for one party than a neutral PSD suggests
that the excess homogeneity is not geographic in origin.

\subsection{Graph-Partitioning with Partisan Edge Weights}

Prior work on edge-weighted redistricting has used edge weights to encode
racial minority clustering (VRA-aligned weighting) \citep{dellib2026vra},
geographic contiguity (standard weighting), and county preservation (county
weighting) \citep{dellib2026county}.  The use of partisan vote similarity as
an edge weight is a natural extension: instead of weighting edges by shared
boundary length or demographic similarity, we weight by partisan similarity.

The only prior work we are aware of that explicitly uses partisan similarity
as a redistricting objective is \citet{gurnee2021fairmandering}, who optimise
directly for a specified partisan outcome.  Our approach differs: we
parametrically vary the similarity weight $\alpha$ (equation~\eqref{eq:psd-weight}),
producing a family of maps from neutral (compact) to maximally homogeneous
(partisan similarity), without targeting any specific partisan outcome.  The
resulting seat distribution is a consequence of geographic sorting, not a
design target.
